\documentclass[11pt,a4paper]{article}

\usepackage[a4paper,top=20mm,bottom=20mm,left=19mm,right=19mm,headheight=15pt]{geometry}
\usepackage{fontspec}
\setsansfont{Open Sans}
\renewcommand{\familydefault}{\sfdefault}
\usepackage{microtype}
\usepackage{xcolor}
\usepackage{graphicx}
\usepackage{array,booktabs,tabularx}
\usepackage{amssymb}
\usepackage{enumitem}
\usepackage{fancyhdr}
\usepackage{titlesec}
\usepackage{tcolorbox}
\tcbuselibrary{breakable,skins}
\usepackage{hyperref}

\definecolor{branddark}{HTML}{02050C}
\definecolor{brandnavy}{HTML}{071426}
\definecolor{brandblue}{HTML}{0066FF}
\definecolor{brandcyan}{HTML}{33AAFF}
\definecolor{brandmint}{HTML}{2DD4BF}
\definecolor{brandviolet}{HTML}{A78BFA}
\definecolor{brandgold}{HTML}{FBBF24}
\definecolor{textdark}{HTML}{172033}
\definecolor{textmuted}{HTML}{53637B}
\definecolor{line}{HTML}{DCE5F0}
\definecolor{paperblue}{HTML}{F2F7FF}

\hypersetup{
  colorlinks=true,
  linkcolor=brandblue,
  urlcolor=brandblue,
  pdftitle={IA na Prática para Empresas},
  pdfauthor={Brimajor Sistemas Embarcados e IA},
  pdfsubject={Capacitação empresarial em IA, engenharia de contexto, processos, agentes e squads}
}

\pagestyle{fancy}
\fancyhf{}
\fancyhead[L]{\small\color{textmuted}BRIMAJOR · CAPACITAÇÃO EMPRESARIAL EM IA}
\fancyhead[R]{\small\color{textmuted}\leftmark}
\fancyfoot[L]{\small\color{textmuted}IA na Prática para Empresas}
\fancyfoot[R]{\small\color{textmuted}\thepage}
\renewcommand{\headrulewidth}{0.3pt}
\renewcommand{\footrulewidth}{0pt}

\titleformat{\section}{\color{branddark}\bfseries\fontsize{24}{28}\selectfont}{\thesection}{0.7em}{}
\titleformat{\subsection}{\color{brandblue}\bfseries\fontsize{15}{18}\selectfont}{\thesubsection}{0.6em}{}
\titlespacing*{\section}{0pt}{0pt}{12pt}
\titlespacing*{\subsection}{0pt}{14pt}{7pt}
\setlength{\parindent}{0pt}
\setlength{\parskip}{6pt}
\setlist[itemize]{leftmargin=5mm,itemsep=3pt,topsep=4pt}
\setlist[enumerate]{leftmargin=6mm,itemsep=4pt,topsep=4pt}
\raggedbottom

\newtcolorbox{destaque}[1][]{
  enhanced,breakable,colback=paperblue,colframe=brandblue!45,
  boxrule=0.7pt,arc=2mm,left=4mm,right=4mm,top=3mm,bottom=3mm,#1
}
\newtcolorbox{atividade}[1][]{
  enhanced,breakable,colback=white,colframe=brandmint,
  boxrule=1pt,arc=2mm,left=4mm,right=4mm,top=3mm,bottom=3mm,
  title={ATIVIDADE PRÁTICA},fonttitle=\bfseries,coltitle=branddark,#1
}
\newtcolorbox{alerta}[1][]{
  enhanced,breakable,colback=brandgold!10,colframe=brandgold!80!black,
  boxrule=0.8pt,arc=2mm,left=4mm,right=4mm,top=3mm,bottom=3mm,#1
}
\newcommand{\linha}{\par\noindent\textcolor{line}{\rule{\linewidth}{0.5pt}}\par}
\newcommand{\campo}[1]{\textbf{#1}\par\vspace{1mm}\textcolor{line}{\rule{\linewidth}{0.6pt}}\par\vspace{2mm}}
\newcommand{\nivel}[2]{\textcolor{#1}{\rule{3mm}{3mm}}\hspace{2mm}\textbf{#2}}

\begin{document}

\begin{titlepage}
  \pagecolor{branddark}\color{white}
  \vspace*{3mm}
  \includegraphics[width=65mm]{../../public/marca/brimajor-logo.png}
  \vfill
  {\color{brandcyan}\rule{24mm}{2pt}}\\[8mm]
  {\fontsize{34}{39}\selectfont\bfseries IA na Prática\\para Empresas}\par
  \vspace{7mm}
  {\fontsize{16}{22}\selectfont\color{white!78}
    Um guia direto sobre engenharia de contexto, uso seguro, melhoria de processos, agentes e squads — sem exigir programação.
  }\par
  \vfill
  \begin{tcolorbox}[colback=brandnavy,colframe=brandcyan!45!branddark,boxrule=0.6pt,arc=2mm,left=5mm,right=5mm,top=4mm,bottom=4mm]
    \color{white}
    \textbf{Público}\quad Empresas começando a usar IA ou prontas para tornar o uso mais consistente.\\[2mm]
    \textbf{Formato}\quad Mini-módulos, exemplos e folhas de trabalho para aplicar em equipe.
  \end{tcolorbox}
  \vspace{8mm}
  {\small\color{white!60}Versão 1.0 · Novembro de 2026 · Manaus, AM}
\end{titlepage}
\nopagecolor\color{textdark}

\section*{Como usar esta apostila}
\addcontentsline{toc}{section}{Como usar esta apostila}

Este material foi feito para uma conversa prática com equipes, não para uma aula de programação. A proposta é sair de ideias abstratas e chegar a um primeiro uso de IA que tenha objetivo, contexto, revisão humana e medida de resultado.

\begin{destaque}
\textbf{Ao final, a equipe deverá conseguir:}
\begin{itemize}
  \item reconhecer tarefas em que a IA ajuda — e aquelas em que não deve decidir sozinha;
  \item montar um contexto completo, reutilizável e verificável;
  \item escolher um processo pequeno para um piloto de 30 dias;
  \item definir regras mínimas de segurança, qualidade e responsabilidade;
  \item distinguir copilotos, agentes e squads de agentes.
\end{itemize}
\end{destaque}

\subsection*{Percurso sugerido}

\begin{tabularx}{\linewidth}{>{\bfseries}p{16mm} >{\bfseries}p{40mm} X p{18mm}}
\toprule
Módulo & Tema & Entrega prática & Tempo \\
\midrule
1 & IA sem mistério & Inventário de tarefas & 35 min \\
2 & Engenharia de contexto & Cartão de contexto & 55 min \\
3 & Uso seguro & Checklist de revisão & 35 min \\
4 & IA nos processos & Canvas de piloto & 50 min \\
5 & Agentes e squads & Mapa de papéis e limites & 35 min \\
\bottomrule
\end{tabularx}

\begin{alerta}
\textbf{Regra de ouro:} a IA pode preparar, resumir, comparar, classificar e sugerir. A responsabilidade por decisões relevantes permanece com pessoas claramente identificadas.
\end{alerta}

\newpage
\section{Mapa de maturidade}
\markboth{Mapa de maturidade}{Mapa de maturidade}

Não existe uma única jornada de IA para todas as empresas. O ponto de partida depende do conhecimento da equipe, da qualidade do processo e do risco da tarefa. Use o mapa abaixo para escolher a trilha adequada.

\begin{tcolorbox}[colback=brandcyan!8,colframe=brandcyan,title={\nivel{brandcyan}{1 · COMEÇAR COM IA}},fonttitle=\bfseries]
\textbf{Para quem:} equipes com zero ou pouco conhecimento.\\
\textbf{Foco:} entender capacidades, limites, segurança e experimentar tarefas simples.\\
\textbf{Sinal de avanço:} a equipe sabe conferir a resposta e não envia dados sensíveis sem autorização.
\end{tcolorbox}

\begin{tcolorbox}[colback=brandmint!8,colframe=brandmint,title={\nivel{brandmint}{2 · ENGENHARIA DE CONTEXTO}},fonttitle=\bfseries]
\textbf{Para quem:} pessoas que já usam ferramentas de IA, mas obtêm resultados instáveis.\\
\textbf{Foco:} combinar objetivo, cenário, referências, restrições, formato e critérios de qualidade.\\
\textbf{Sinal de avanço:} bons padrões são documentados e reutilizados pela equipe.
\end{tcolorbox}

\begin{tcolorbox}[colback=brandviolet!8,colframe=brandviolet,title={\nivel{brandviolet}{3 · IA NOS PROCESSOS}},fonttitle=\bfseries]
\textbf{Para quem:} empresas que já usam IA no dia a dia.\\
\textbf{Foco:} transformar usos individuais em rotinas supervisionadas, mensuráveis e melhoráveis.\\
\textbf{Sinal de avanço:} existe um dono do processo, indicadores e tratamento de exceções.
\end{tcolorbox}

\begin{tcolorbox}[colback=brandgold!10,colframe=brandgold!90!black,title={\nivel{brandgold!90!black}{4 · AGENTES E SQUADS}},fonttitle=\bfseries]
\textbf{Para quem:} equipes com processos estáveis, dados organizados e governança mínima.\\
\textbf{Foco:} agentes com ferramentas e squads de agentes especializados, sempre com limites e supervisão.\\
\textbf{Sinal de prontidão:} a empresa consegue explicar quem autoriza, quem revisa e como interromper a operação.
\end{tcolorbox}

\newpage
\section{IA sem mistério}
\markboth{IA sem mistério}{IA sem mistério}

Uma IA generativa produz uma resposta provável a partir do pedido e do contexto disponíveis. Ela não “sabe” a verdade como uma pessoa e não entende a empresa por padrão. Seu resultado melhora quando recebe boas referências, regras claras e exemplos relevantes.

\begin{destaque}
\textbf{Um modelo de IA é excelente para gerar possibilidades; não é, sozinho, uma fonte de verdade.}
\end{destaque}

\subsection{Onde ela costuma ajudar}

\begin{tabularx}{\linewidth}{>{\bfseries\color{brandblue}}p{34mm} X}
Preparar & rascunhos, pautas, roteiros, perguntas e estruturas iniciais. \\
Resumir & documentos longos, reuniões e históricos — desde que o conteúdo seja fornecido. \\
Transformar & adaptar tom, formato, público, idioma ou nível de detalhe. \\
Comparar & organizar alternativas segundo critérios definidos por uma pessoa. \\
Classificar & separar itens em categorias conhecidas e sinalizar exceções. \\
Extrair & localizar campos, compromissos, riscos ou temas em um conjunto de textos. \\
\end{tabularx}

\subsection{Onde exige mais cuidado}

\begin{itemize}
  \item decisões jurídicas, médicas, financeiras, trabalhistas ou de segurança;
  \item fatos que precisam estar atualizados ou acompanhados de fonte confiável;
  \item avaliações de pessoas, concessão de acesso ou liberação de dinheiro;
  \item tarefas com dados pessoais, segredos comerciais ou credenciais;
  \item situações em que um erro não pode ser detectado antes de causar dano.
\end{itemize}

\begin{alerta}
Se o custo de um erro for alto, a solução não é “um prompt melhor”. É preciso reduzir o escopo, usar fontes controladas, exigir aprovação humana ou não automatizar.
\end{alerta}

\newpage
\subsection{Escolha tarefas, não modismos}

Uma boa oportunidade tem volume suficiente para justificar o esforço, padrão reconhecível, entrada disponível e uma pessoa capaz de revisar a saída. Evite começar pelo processo mais complexo ou mais crítico da empresa.

\begin{atividade}
Liste tarefas recorrentes da equipe e faça uma triagem inicial.

\begin{tabularx}{\linewidth}{p{43mm} c c c X}
\toprule
Tarefa & Repete? & Revisável? & Baixo risco? & Primeiro teste \\
\midrule
\rule{0pt}{7mm} & $\square$ & $\square$ & $\square$ & \\
\rule{0pt}{7mm} & $\square$ & $\square$ & $\square$ & \\
\rule{0pt}{7mm} & $\square$ & $\square$ & $\square$ & \\
\rule{0pt}{7mm} & $\square$ & $\square$ & $\square$ & \\
\bottomrule
\end{tabularx}
\end{atividade}

\subsection*{Teste rápido de adequação}

Marque um ponto para cada “sim”. Comece por tarefas com quatro ou cinco pontos.

\begin{enumerate}
  \item A entrada está disponível em formato digital?
  \item É possível descrever o que caracteriza uma boa saída?
  \item Uma pessoa consegue revisar antes do uso?
  \item O erro pode ser corrigido sem dano relevante?
  \item A tarefa acontece com frequência suficiente para gerar aprendizado?
\end{enumerate}

\begin{destaque}
\textbf{Exemplos de bons primeiros pilotos:} resumir atas; adaptar uma comunicação para públicos diferentes; preparar perguntas para uma reunião; classificar solicitações para triagem; comparar propostas segundo critérios definidos.
\end{destaque}

\newpage
\section{Engenharia de contexto}
\markboth{Engenharia de contexto}{Engenharia de contexto}

Um comando isolado diz o que fazer. A engenharia de contexto organiza tudo o que a IA precisa para fazer bem: objetivo, público, cenário, referências, limites, formato e critérios de verificação.

\begin{center}
\begin{tcolorbox}[width=.92\linewidth,colback=branddark,colframe=branddark,arc=2mm]
\centering\color{white}\large\bfseries
Resultado útil = tarefa clara + contexto confiável + limites + revisão
\end{tcolorbox}
\end{center}

\subsection{As sete camadas do contexto}

\begin{tabularx}{\linewidth}{>{\bfseries\color{brandblue}}p{36mm} X}
1. Objetivo & O que precisa mudar ou ficar pronto ao final? \\
2. Público & Quem vai usar ou receber a saída? O que essa pessoa já sabe? \\
3. Cenário & Qual é a situação, o momento e o processo em que a tarefa acontece? \\
4. Referências & Quais documentos, exemplos, dados e fontes devem orientar a resposta? \\
5. Restrições & O que não pode acontecer? Que política, prazo, tom ou limite se aplica? \\
6. Formato & Como a saída deve ser entregue: tabela, e-mail, lista, plano, campos? \\
7. Verificação & Que checklist permite conferir qualidade, fontes, lacunas e riscos? \\
\end{tabularx}

\begin{alerta}
Contexto não é quantidade de texto. É informação relevante, atual, autorizada e organizada. Excesso de material contraditório também piora a resposta.
\end{alerta}

\newpage
\subsection{Cartão de contexto reutilizável}

Use um cartão por tarefa recorrente. Registre o dono, as fontes e a data de revisão para que a equipe não dependa da memória de uma única pessoa.

\begin{atividade}[title={CARTÃO DE CONTEXTO}]
\campo{Tarefa e resultado esperado}
\campo{Público e situação de uso}
\campo{Entradas e referências autorizadas}
\campo{Regras, limites e itens proibidos}
\campo{Formato da saída}
\campo{Critérios de qualidade e fatos que devem ser conferidos}
\campo{Responsável pela revisão humana}
\textbf{Dono do cartão:} \rule{45mm}{0.5pt}\hfill
\textbf{Revisar em:} \rule{30mm}{0.5pt}
\end{atividade}

\newpage
\subsection{Do pedido vago ao trabalho verificável}

\begin{tcolorbox}[colback=red!3,colframe=red!45!black,title={ANTES · pedido vago},fonttitle=\bfseries]
“Faça um e-mail para nossos clientes sobre a mudança.”
\end{tcolorbox}

\begin{tcolorbox}[colback=brandmint!6,colframe=brandmint,title={DEPOIS · contexto organizado},fonttitle=\bfseries]
\textbf{Objetivo:} avisar clientes ativos sobre a mudança do horário de atendimento a partir de 1º de dezembro.\\
\textbf{Público:} clientes empresariais que já usam nosso suporte.\\
\textbf{Referência:} comunicado interno aprovado abaixo. Use somente os fatos presentes nele.\\
\textbf{Tom:} direto, cordial e sem linguagem técnica.\\
\textbf{Formato:} assunto com até 55 caracteres e corpo com até 140 palavras.\\
\textbf{Obrigatório:} informar novo horário, data de início e canal para dúvidas.\\
\textbf{Verificação:} ao final, liste qualquer informação ausente em “Pendências”; não invente.
\end{tcolorbox}

\subsection*{Por que o segundo funciona melhor?}

Ele reduz interpretações, traz uma fonte, define o que não inventar e facilita a revisão. A qualidade deixa de depender apenas de “saber conversar com a IA” e passa a depender de um padrão que a equipe pode melhorar.

\begin{destaque}
\textbf{Prática recomendada:} guarde cartões, exemplos aprovados e checklists em uma biblioteca de contexto com responsáveis, versão e validade.
\end{destaque}

\newpage
\section{Uso seguro e responsável}
\markboth{Uso seguro}{Uso seguro}

Segurança começa antes de abrir a ferramenta. A empresa precisa saber quais serviços estão autorizados, que tipo de dado pode ser enviado e quem responde pela revisão.

\subsection{Semáforo de dados}

\begin{tabularx}{\linewidth}{>{\bfseries}p{28mm} X X}
\toprule
Faixa & Exemplos & Conduta \\
\midrule
\textcolor{green!50!black}{VERDE} & conteúdo público, material de treinamento, dados fictícios & pode usar em ferramenta autorizada, mantendo revisão. \\
\textcolor{brandgold!80!black}{AMARELO} & informação interna, contratos, indicadores, dados de cliente já anonimizados & usar somente com autorização, controles e finalidade definida. \\
\textcolor{red!70!black}{VERMELHO} & senhas, chaves, documentos pessoais, dados de saúde, segredo industrial, informação regulada & não enviar sem solução aprovada, base legal e proteção compatível. \\
\bottomrule
\end{tabularx}

\subsection{Quatro riscos frequentes}

\begin{enumerate}
  \item \textbf{Invenção convincente:} a resposta parece segura, mas contém fatos inexistentes.
  \item \textbf{Exposição:} dados foram enviados a uma ferramenta sem autorização ou configuração adequada.
  \item \textbf{Viés:} o modelo reproduz padrões injustos presentes nas referências ou na sociedade.
  \item \textbf{Automação silenciosa:} ninguém sabe que a IA participou e não existe responsável pela decisão.
\end{enumerate}

\begin{alerta}
Nunca compartilhe credenciais, segredos ou dados pessoais apenas porque a tarefa parece urgente. Anonimizar não é só apagar o nome: combinações de cargo, data, local e caso podem reidentificar uma pessoa.
\end{alerta}

\newpage
\subsection{Fluxo mínimo de verificação}

\begin{center}
\begin{tabularx}{.96\linewidth}{>{\centering\arraybackslash\bfseries}X >{\centering\arraybackslash}p{10mm} >{\centering\arraybackslash\bfseries}X >{\centering\arraybackslash}p{10mm} >{\centering\arraybackslash\bfseries}X}
PEDIDO + CONTEXTO & $\rightarrow$ & RASCUNHO DA IA & $\rightarrow$ & REVISÃO HUMANA \\
\end{tabularx}
\vspace{3mm}
\begin{tabularx}{.96\linewidth}{>{\centering\arraybackslash\bfseries}X >{\centering\arraybackslash}p{10mm} >{\centering\arraybackslash\bfseries}X >{\centering\arraybackslash}p{10mm} >{\centering\arraybackslash\bfseries}X}
REGISTRO & $\leftarrow$ & USO APROVADO & $\leftarrow$ & CORREÇÃO \\
\end{tabularx}
\end{center}

\begin{atividade}[title={CHECKLIST ANTES DE USAR A SAÍDA}]
\begin{itemize}[label=$\square$]
  \item Os fatos principais foram conferidos na fonte original?
  \item Números, datas, nomes e links estão corretos?
  \item O texto atende ao objetivo e ao público definidos?
  \item Alguma informação foi inventada, extrapolada ou omitida?
  \item A saída expõe dado pessoal, confidencial ou desnecessário?
  \item Uma pessoa identificada aprovou o uso?
  \item Vale registrar o erro ou a melhoria na biblioteca de contexto?
\end{itemize}
\end{atividade}

\begin{destaque}
\textbf{Política mínima de uma página:} ferramentas autorizadas; dados permitidos e proibidos; usos que exigem aprovação; responsável por incidentes; canal para dúvidas; periodicidade de revisão.
\end{destaque}

\newpage
\section{IA nos processos}
\markboth{IA nos processos}{IA nos processos}

O uso intermediário começa quando a empresa deixa de depender de iniciativas isoladas e passa a desenhar uma rotina. A IA ocupa etapas específicas; as pessoas continuam responsáveis por objetivos, exceções e decisões.

\subsection{Mapeie antes de automatizar}

\begin{enumerate}
  \item \textbf{Gatilho:} o que inicia o processo?
  \item \textbf{Entrada:} que informação chega e em que qualidade?
  \item \textbf{Etapas:} onde existem leitura, transformação, comparação ou classificação?
  \item \textbf{Decisões:} quais escolhas exigem julgamento ou autorização?
  \item \textbf{Exceções:} o que sai do padrão e para quem deve ser encaminhado?
  \item \textbf{Saída:} o que precisa ficar pronto e como se comprova a qualidade?
\end{enumerate}

\begin{destaque}
\textbf{Divisão saudável do trabalho}\par
\vspace{2mm}
\begin{tabularx}{\linewidth}{>{\bfseries\color{brandblue}}p{30mm} X}
IA & prepara, extrai, resume, propõe, classifica e sinaliza. \\
Pessoa & define, autoriza, investiga exceções, corrige e assume a decisão. \\
Sistema & controla acesso, registra versões, aplica regras e guarda evidências. \\
\end{tabularx}
\end{destaque}

\subsection{Critério de um bom primeiro piloto}

Escolha um processo com dono definido, baixo risco, dados disponíveis, frequência semanal e uma medida de sucesso simples. Mantenha a alternativa manual durante o teste.

\newpage
\subsection{Canvas do piloto}

\begin{atividade}[title={PILOTO DE IA · UMA PÁGINA}]
\campo{Problema concreto e pessoas afetadas}
\campo{Etapa em que a IA ajudará — e o que ficará fora}
\campo{Entradas, fontes autorizadas e cartão de contexto}
\campo{Saída esperada e critérios de qualidade}
\campo{Revisor, aprovador e tratamento das exceções}
\campo{Indicador atual e meta para 30 dias}
\campo{Riscos, forma de interrupção e alternativa manual}
\end{atividade}

\subsection*{Exemplo de hipótese}

“Se usarmos um cartão de contexto para preparar o primeiro rascunho das atas, reduziremos o tempo médio de 45 para 20 minutos, mantendo pelo menos 95\% dos compromissos corretos após revisão humana.”

Uma boa hipótese contém tarefa, mudança, medida, prazo e condição de qualidade. “Usar IA para ser mais produtivo” não permite saber se o piloto funcionou.

\newpage
\subsection{Da experiência individual ao padrão da equipe}

Depois do piloto, transforme o aprendizado em ativos simples:

\begin{itemize}
  \item cartão de contexto aprovado e exemplos de boa saída;
  \item checklist de revisão e lista das falhas mais frequentes;
  \item decisão sobre dados, ferramenta e acesso;
  \item procedimento para exceções e escalonamento;
  \item métrica acompanhada e responsável pelo ciclo de melhoria.
\end{itemize}

\subsection{Indicadores úteis}

\begin{tabularx}{\linewidth}{>{\bfseries\color{brandblue}}p{35mm} X X}
Resultado & O que medir & Cuidado \\
\midrule
Tempo & minutos por tarefa, fila, prazo de resposta & não confundir velocidade com qualidade. \\
Qualidade & taxa de aprovação, correções, completude & amostra e critério precisam ser estáveis. \\
Adoção & pessoas ativas, frequência, processos cobertos & uso alto pode esconder retrabalho. \\
Risco & incidentes, dados bloqueados, exceções & incentive o registro, não a ocultação. \\
Valor & custo evitado, receita apoiada, capacidade liberada & compare com uma linha de base realista. \\
\end{tabularx}

\begin{alerta}
Não premie apenas “quantidade de uso”. A meta é melhorar um resultado de negócio com qualidade e controle, não aumentar o número de mensagens enviadas à IA.
\end{alerta}

\newpage
\section{Copilotos, agentes e squads}
\markboth{Agentes e squads}{Agentes e squads}

Estas formas de trabalho não são degraus obrigatórios para toda empresa. Quanto maior a autonomia, maior a necessidade de processo estável, permissões restritas, testes, observabilidade e supervisão.

\begin{tabularx}{\linewidth}{>{\bfseries}p{30mm} X X}
\toprule
Modelo & Como atua & Quando considerar \\
\midrule
Assistente & responde a um pedido por vez & exploração e tarefas pontuais. \\
Copiloto & apoia etapas de um fluxo conduzido por uma pessoa & rotina conhecida com revisão constante. \\
Agente & recebe um objetivo, usa ferramentas e executa várias etapas dentro de limites & processo estável, permissões controladas e resultado verificável. \\
Squad de agentes & distribui o fluxo entre papéis especializados e revisores & trabalho complexo que realmente se beneficia da especialização. \\
\bottomrule
\end{tabularx}

\subsection{Anatomia de um agente supervisionado}

\begin{itemize}
  \item \textbf{Objetivo:} tarefa delimitada e condição clara de conclusão.
  \item \textbf{Contexto:} fontes autorizadas, exemplos, regras e critérios.
  \item \textbf{Ferramentas:} acesso mínimo necessário, com permissões separadas.
  \item \textbf{Estado:} registro do que já foi feito, sem memória ilimitada ou secreta.
  \item \textbf{Limites:} orçamento, prazo, número de tentativas e ações proibidas.
  \item \textbf{Supervisão:} pontos em que uma pessoa aprova, corrige ou interrompe.
  \item \textbf{Evidência:} histórico de fontes, ações, versões e decisões.
\end{itemize}

\newpage
\subsection{Como pensar uma squad para a empresa}

Uma squad de agentes organiza papéis, não apenas várias IAs conversando. Cada papel precisa de entrada, saída, ferramentas, critérios de passagem e responsável humano.

\begin{center}
\begin{tabularx}{.96\linewidth}{>{\centering\arraybackslash\bfseries}X >{\centering\arraybackslash}p{8mm} >{\centering\arraybackslash\bfseries}X >{\centering\arraybackslash}p{8mm} >{\centering\arraybackslash\bfseries}X}
PLANEJADOR & $\rightarrow$ & ESPECIALISTA & $\rightarrow$ & REVISOR \\
\small decompõe e define critérios & & \small produz com ferramentas limitadas & & \small verifica contra evidências \\
\end{tabularx}
\vspace{5mm}
\fcolorbox{brandblue}{paperblue}{\parbox{.86\linewidth}{\centering\textbf{PESSOA RESPONSÁVEL}\quad autoriza mudanças relevantes, resolve exceções e responde pelo resultado}}
\end{center}

\subsection*{Perguntas antes de montar}

\begin{enumerate}
  \item Uma única etapa assistida já resolve o problema?
  \item O processo e os critérios estão documentados?
  \item Quais ações podem apenas ser propostas e quais podem ser executadas?
  \item Que agente pode acessar qual fonte ou ferramenta?
  \item Onde a aprovação humana é obrigatória?
  \item Como detectar, registrar, interromper e recuperar uma falha?
\end{enumerate}

\begin{alerta}
Não dê a um agente autonomia ampla só porque ele “pareceu funcionar” em uma demonstração. Teste casos normais, exceções, entradas maliciosas, indisponibilidade e limites de permissão.
\end{alerta}

\newpage
\section{Plano de ação em 30 dias}
\markboth{Plano de 30 dias}{Plano de 30 dias}

\begin{tabularx}{\linewidth}{>{\bfseries\color{brandblue}}p{28mm} X p{38mm}}
\toprule
Período & Ação & Evidência de saída \\
\midrule
Dias 1–5 & formar um pequeno grupo, mapear tarefas e definir regras mínimas de dados & inventário priorizado e política provisória. \\
Dias 6–10 & escolher uma tarefa, medir a linha de base e montar o cartão de contexto & hipótese e exemplo aprovado. \\
Dias 11–20 & executar o piloto em amostra controlada, sempre com revisão humana & registro de tempo, correções e falhas. \\
Dias 21–25 & comparar com a linha de base, ouvir usuários e corrigir o padrão & relatório curto e versão 2 do cartão. \\
Dias 26–30 & decidir: encerrar, repetir ou ampliar com novos controles & decisão registrada, dono e próximo marco. \\
\bottomrule
\end{tabularx}

\begin{atividade}[title={COMPROMISSO DA EQUIPE}]
\campo{Tarefa escolhida para o primeiro piloto}
\campo{Resultado que queremos melhorar}
\campo{Pessoa dona do processo e pessoa revisora}
\campo{Primeira entrega e data}
\campo{Como saberemos se vale continuar}
\end{atividade}

\begin{destaque}
\textbf{Decisão ao final:}
\begin{itemize}
  \item \textbf{encerrar} quando o valor não compensa o risco ou o esforço;
  \item \textbf{repetir} quando há potencial, mas o contexto ou o processo ainda são instáveis;
  \item \textbf{ampliar} somente quando qualidade, responsabilidade e controles estão demonstrados.
\end{itemize}
\end{destaque}

\newpage
\section*{Resumo de bolso}
\addcontentsline{toc}{section}{Resumo de bolso}
\markboth{Resumo de bolso}{Resumo de bolso}

\begin{tcolorbox}[colback=branddark,colframe=branddark,arc=2mm,left=6mm,right=6mm,top=5mm,bottom=5mm]
\color{white}
\begin{enumerate}
  \item Comece por uma tarefa pequena, recorrente e revisável.
  \item Não trate a IA como fonte de verdade.
  \item Organize objetivo, público, cenário, referências, limites, formato e verificação.
  \item Use somente ferramentas e dados autorizados.
  \item Mantenha uma pessoa responsável por aprovar o resultado.
  \item Meça qualidade e impacto contra uma linha de base.
  \item Registre falhas e transforme aprendizados em contexto reutilizável.
  \item Só avance para agentes quando o processo já for estável.
  \item Em squads, defina papéis, permissões, passagens e revisão humana.
  \item Preserve sempre uma forma de interromper e recuperar a operação.
\end{enumerate}
\end{tcolorbox}

\subsection*{Glossário essencial}

\begin{tabularx}{\linewidth}{>{\bfseries\color{brandblue}}p{38mm} X}
Contexto & informação organizada que orienta a IA: fontes, regras, exemplos, cenário e critérios. \\
Alucinação & conteúdo inventado ou incorreto apresentado de forma convincente. \\
Copiloto & IA que apoia uma pessoa dentro de uma rotina. \\
Agente & sistema que planeja e executa etapas com ferramentas dentro de limites definidos. \\
Squad de agentes & conjunto de papéis especializados que colaboram em um fluxo supervisionado. \\
Supervisão humana & pontos explícitos de revisão, aprovação, correção e interrupção. \\
\end{tabularx}

\vfill
\begin{center}
\includegraphics[width=45mm]{../../public/marca/brimajor-icone.png}\\[3mm]
{\bfseries\color{brandblue}BRIMAJOR SISTEMAS EMBARCADOS \& IA}\\
{\small\color{textmuted}Capacitação prática para transformar IA em trabalho bem definido, seguro e mensurável.}
\end{center}

\end{document}
